(1) 约束给 \(\dot z=2\alpha r\dot r\)，轨道旋转使速度平方为 \((1+4\alpha^2r^2)\dot r^2+\Omega^2r^2\)。因此
\[
\mathcal L=\frac m2(1+4\alpha^2r^2)\dot r^2+\frac m2\Omega^2r^2-mg\alpha r^2.
\]
有效势 \(U_{\rm eff}=m(g\alpha-\Omega^2/2)r^2\)，虽看似谐振，动能的有效质量仍依赖 r。

(2) 平衡条件 \(U'_{\rm eff}=m(2g\alpha-\Omega^2)r=0\)，因此只要 \(\Omega^2\ne2g\alpha\)，唯一平衡是 r=0；临界 \(\Omega_c=\sqrt{2g\alpha}\) 时所有 r 都满足静态势能驻值，是理想模型中的连续中性平衡，而不是有限位置的两个新稳定分支。

(3) Euler–Lagrange 方程为
\[
(1+4\alpha^2r^2)\ddot r+4\alpha^2r\dot r^2+(2g\alpha-\Omega^2)r=0.
\]
底点附近 \(\ddot r+(2g\alpha-\Omega^2)r=0\)，所以 \(\omega^2=2g\alpha-\Omega^2\)。低于临界稳定，高于临界线性不稳定，临界频率为零。

(4) 令 \(\Omega=0\) 后 \(U_{\rm eff}=mg\alpha r^2\) 正定，\(\omega=\sqrt{2g\alpha}\)。直接用重力沿轨道的切向投影在 \(r\approx 0\) 得 \(\ddot r=-2g\alpha r\)，与变分方程相符。
